% Part: computability
% Chapter: recursive-functions
% Section: sequences

\documentclass[../../../include/open-logic-section]{subfiles}

\begin{document}

\olfileid{cmp}{rec}{seq}
\olsection{Rerangken}

Himpunan fungsi rekursif primitif iku pancen kukuh.
Nanging kita bisa nindakake luwih akeh maneh sawisé
duwé cara kang nyukupi kanggo ngolah \emph{rerangken}.
Kita bakal nggandhengake rerangken winates saka wilangan
asli karo wilangan asli kanthi cara iki: rerangken
$\langle a_0, a_1, a_2, \dots, a_k \rangle$ cocog karo
wilangan
\[
p_0^{a_0+1} \cdot p_1^{a_1+1} \cdot p_2^{a_2+1} \cdot \dots \cdot
p_k^{a_k+1}.
\]
Kita nambah siji ing saben pangkat supaya, upamane,
rerangken $\langle 2, 7, 3\rangle$ lan
$\langle 2, 7, 3, 0, 0 \rangle$ nduwé kode numerik
kang béda. Kita bisa nggunakake $0$ utawa~$1$ kanggo
ngode rerangken kosong; supaya cetha, $\emptyseq$
nglambangake~$0$.

Pangodean rerangken iki bisa lumaku amarga Teorema
Dhasar Aritmetika: saben wilangan asli $n \ge 2$ bisa
ditulis kanthi siji cara waé ing wangun
\[
n = p_0^{a_0} \cdot p_1^{a_1} \cdot \dots \cdot p_k^{a_k}
\]
kanthi $a_k \ge 1$. Iki njamin yen pemetaan
$\tuple{}(a_0, \dots, a_k) = \tuple{a_0, \dots, a_k}$
iku injektif: rerangken kang béda dipetakake menyang
wilangan kang béda; saben wilangan cocog karo paling
akeh siji rerangken.

Saiki kita bakal nuduhake yen operasi nemtokake dawa
rerangken, nemtokake unsur kaping~$i$, nambahi unsur
ing pucuk rerangken, lan nyambungake rong rerangken,
kabeh rekursif primitif.

\begin{prop}
  Fungsi $\len{s}$, kang ngasilake dawa rerangken
  $s$, iku rekursif primitif.
\end{prop}

\begin{proof}
  Upamana $R(i, s)$ iku relasi kang ditegesi kanthi
  \[
  R(i, s) \text{ yen lan mung yen }
  p_i \mid s \land p_{i+1} \nmid s.
  \]
  $R$ cetha rekursif primitif. Yen $s$ dadi kode
  rerangken kang ora kosong, yaiku
  \[
  s = p_0^{a_0+1} \cdot \dots \cdot p_{k}^{a_{k}+1},
  \]
  $R(i,s)$ bener yen $p_i$ iku prima paling gedhe
  kang nyukupi $p_i \mid s$, yaiku $i = k$.
  Dadi dawa $s$ iku $i+1$ yen lan mung yen
  $p_i$ prima paling gedhe kang mbagi~$s$,
  mula kita bisa netepake
  \[
  \len{s} =
  \begin{cases}
    0 & \text{yen $s = 0$ utawa $s = 1$} \\
    1 + \bmin{i < s}{R(i, s)} & \text{yen ora}
  \end{cases}
  \]
  Ing kéné kita bisa nggunakake minimisasi mawa wates,
  amarga mung ana siji $i$ kang nyukupi $R(i,s)$
  nalika $s$ kode rerangken, lan yen $i$ ana,
  nilaine luwih cilik tinimbang~$s$.
\end{proof}

\begin{prop}
  Fungsi $\fn{append}(s,a)$, kang ngasilake rerangken
  sawisé $a$ ditambahake ing pucuk rerangken $s$,
  iku rekursif primitif.
\end{prop}

\begin{proof}
  $\fn{append}$ bisa ditegesi kanthi:
  \[
  \fn{append}(s,a) =
  \begin{cases}
    2^{a+1} & \text{yen $s = 0$ utawa $s = 1$} \\
    s \cdot p_{\len{s}}^{a+1} & \text{yen ora.}
  \end{cases}
  \]
\end{proof}

\begin{prop}
  Fungsi $\fn{element}(s,i)$, kang ngasilake unsur
  kaping~$i$ saka $s$ (unsur wiwitan diarani
  unsur kaping~$0$), utawa $0$ yen $i$ luwih
  gedhe tinimbang utawa padha karo dawa $s$,
  iku rekursif primitif.
\end{prop}

\begin{proof}
Gatekna yen $a$ iku unsur kaping~$i$ saka~$s$
yen lan mung yen $p_i^{a+1}$ iku pangkat~$p_i$
paling gedhe kang mbagi~$s$, yaiku
$p_i^{a+1} \mid s$ nanging
$p_i^{a+2} \nmid s$. Dadi:
  \[
  \fn{element}(s,i) =
  \begin{cases}
    0 & \mbox{yen $i \geq \len{s}$} \\
    \bmin{a < s}{(p_i^{a+2} \nmid s)} & \text{yen ora.}
  \end{cases}
  \]
\end{proof}

Tinimbang nggunakake jeneng resmi fungsi-fungsi
kang wis ditegesi mau, kita ngenalake notasi kang
luwih ringkes. Kita bakal nulis $(s)_i$ kanggo
$\fn{element}(s,i)$, lan
$\tuple{s_0, \dots, s_k}$ kanggo ngringkes
\[
\fn{append}(\fn{append}(\dots \fn{append}(\emptyseq,s_0)
\dots),s_k).
\]
Gatekna yen dawa $s$ yaiku~$k$, unsur-unsur $s$ yaiku
$(s)_0$, \dots,~$(s)_{k-1}$.

\begin{prop}
Fungsi $\fn{concat}(s,t)$, kang nyambungake rong
rerangken, iku rekursif primitif.
\end{prop}

\begin{proof}
  Kita ngupaya fungsi $\fn{concat}$ kang nyukupi
  \[
    \fn{concat}(\tuple{a_0, \dots, a_k}, \tuple{b_0, \dots, b_l}) =
    \tuple{a_0, \dots, a_k, b_0, \dots, b_l}.
  \]
  Kita nggunakake fungsi pambantu
  $\fn{hconcat}(s,t,n)$ kang nyambungake $n$
  unsur wiwitan saka $t$ menyang~$s$.
  Fungsi iki bisa ditegesi kanthi rekursi
  primitif kaya mangkene:
  \begin{align*}
    \fn{hconcat}(s,t,0) & = s\\
    \fn{hconcat}(s,t,n+1) & = \fn{append}(\fn{hconcat}(s,t,n),(t)_n)
    \intertext{Banjur kita bisa netepake $\fn{concat}$ kanthi}
    \fn{concat}(s,t) & = \fn{hconcat}(s,t,\len{t}).
  \end{align*}
\end{proof}

Kita bakal nulis $s \concat t$ kanggo
$\fn{concat}(s,t)$.

Migunani yen kita bisa menehi wates kanggo kode
numerik rerangken adhedhasar dawane lan unsur
paling gedhene. Upamana $s$ rerangken kang
dawane~$k$, lan saben unsuré kurang saka utawa
padha karo sawenehing wilangan~$x$. Dadi $s$
nduwé paling akeh $k$ faktor prima, saben faktor
ora ngluwihi~$p_{k-1}$, lan pangkate ora
ngluwihi $x+1$ ing faktorisasi prima~$s$.
Kanthi tembung liya, kanggo dawa positif, yen
kita netepake
\[
\fn{sequenceBound}(x,k) = p_{k-1}^{k \cdot (x+1)},
\]
kode numerik rerangken~$s$ kang kasebut mau
ora ngluwihi~$\fn{sequenceBound}(x,k)$.
Yen dawane nol, kita netepake nilai wates iki
dadi siji supaya uga nyakup kode nol kanggo
rerangken kosong.
% OLPL-128: p_{k-1} ora ditegesi kanggo k=0;
% nilai wates siji nutupi kode kosong nol.

Wates rerangken iki menehi cara kanggo netepake
fungsi-fungsi anyar kanthi panggolekan mawa wates.
Umpamane, kita bisa netepake $\fn{concat}$
kanthi panggolekan mawa wates. Kita mung perlu
nulis \emph{spesifikasi} rekursif primitif kanggo
obyek kang digoleki (kode numerik rerangken
gandhengan), lan wates panggolekane. Iki kang
nyukupi:
\begin{align*}
  \fn{concat}(s,t) = {} & \bmin{v \leq \fn{sequenceBound}(s+t,\len{s} +
    \len{t})}{} \\
  & \quad(\len{v} = \len{s} + \len{t} \land {}\\
    & \qquad \bforall{i < \len{s}}{((v)_i = (s)_i)} \land {} \\
      & \qquad \bforall{j < \len{t}}{((v)_{\len{s}+j} = (t)_j)})
\end{align*}
% OLPL-129: Wates sumber nganggo v < B, sanadyan kode bisa padha B.
% OLPL-130: Rong kuantor kudu sejajar supaya unsur t tetep dipriksa
% nalika s kosong.

\begin{prob}
Tuduhna yen ana fungsi rekursif primitif~$\fn{sconcat}(s)$
kang nyukupi
\[
\fn{sconcat}(\tuple{s_0, \dots, s_k}) = s_0 \concat \dots \concat s_k.
\]
\end{prob}

\begin{prob}
Tuduhna yen ana fungsi rekursif primitif~$\fn{tail}(s)$
kang nyukupi
\begin{align*}
  \fn{tail}(\emptyseq) & = 0 \text{ lan}\\
  \fn{tail}(\tuple{s_0, \dots, s_{k}}) & = \tuple{s_1, \dots, s_{k}}.
\end{align*}
\end{prob}

\begin{prop}
  \ollabel{prop:subseq}
  Fungsi $\fn{subseq}(s, i, n)$ kang ngasilake
  rerangken bagean saka $s$ kang dawane~$n$ wiwit
  unsur kaping~$i$, iku rekursif primitif.
\end{prop}

\begin{proof}
  Latihan.
\end{proof}

\begin{prob}
  Buktikna \olref[cmp][rec][seq]{prop:subseq}.
\end{prob}
  
\end{document}
